\documentclass[11pt,a4paper]{article}
\usepackage[T1]{fontenc}
\usepackage[utf8]{inputenc}
\usepackage{amsmath,amssymb,amsthm}
\usepackage[margin=1in]{geometry}
\usepackage[colorlinks=true,linkcolor=blue,urlcolor=blue]{hyperref}
\theoremstyle{plain}
\newtheorem{theorem}{Theorem}
\newtheorem{lemma}[theorem]{Lemma}
\newtheorem{proposition}[theorem]{Proposition}
\newtheorem{corollary}[theorem]{Corollary}
\theoremstyle{definition}
\newtheorem{remark}[theorem]{Remark}
% A quoted conjecture can be a single hundred-term formula whose terms are thirty-digit
% numbers. TeX cannot hyphenate those, so without a little stretch every such quote runs off
% the page; the linked-a-file papers are all of that shape.
\setlength{\emergencystretch}{3em}

\title{The empirical recurrence for OEIS A234180, proved by transfer matrix}
\author{Adrian Perez Fontelles\\ \small Independent researcher}
\date{14 September 2026}
\begin{document}
\maketitle

\begin{abstract}
OEIS A234180 counts $(n+1)\times5$ arrays over $\{0,\dots,5\}$ subject to a condition imposed on
every $2\times2$ subblock, and carries an empirical recurrence of order $60$ contributed
by R.~H.~Hardin. It is true, and it is not an empirical matter at all. The condition is local
to each subblock, so the rows of an admissible array are the vertices of a finite digraph
on $S=7776$ vertices whose edges are the admissible consecutive pairs, and the count is a
number of walks: $a(n)=\mathbf{1}^{\!\top}M^{n}\mathbf{1}$. Writing $q$
for the characteristic polynomial of the conjectured recurrence, the recurrence holds for all
$n$ exactly when $\mathbf{1}^{\!\top}M^{j}q(M)\mathbf{1}=0$ for every $j\ge0$, and the
Cayley--Hamilton theorem bounds that to $j<S$. It is therefore a finite computation in exact
integer arithmetic, and it comes out zero.
\end{abstract}

\noindent\small 2020 Mathematics Subject Classification. 05A15, 05B45, 15B36.\normalsize

\section{The sequence and the conjecture}

OEIS A234180 is ``Number of (n+1)X(4+1) 0..5 arrays with every 2X2 subblock having the absolute values of all six edge and diagonal differences no larger than 1.''. It has offset $1$ and begins
\[
10120, 392664, 15322566, 604124712, 23888919458, 946183426722, 37514052801518, 1488477618052368,\ \dots
\]
The entry carries the following comment:
\begin{quote}\small
Empirical recurrence of order 60 (see link above)

Empirical: a(n) = 149*a(n-1) - 8814*a(n-2) + 256435*a(n-3) - 3393920*a(n-4) + 449245*a(n-5) + 514842621*a(n-6) - 4157999842*a(n-7) - 23261708552*a(n-8) + 407861534972*a(n-9) + 96979298270*a(n-10) - 20363478746809*a(n-11) + 30533472981167*a(n-12) + 660568556988005*a(n-13) - 1408051468806556*a(n-14) - 15383126507976075*a(n-15) + 32902261921526205*a(n-16) + 269077534259492089*a(n-17) - 467962804548993846*a(n-18) - 3572775805160926106*a(n-19) + 4038250887549205858*a(n-20) + 35582810698357101463*a(n-21) - 17005381783517699700*a(n-22) - 259443464221965708471*a(n-23) - 32542572652022730119*a(n-24) + 1341569768136539281807*a(n-25) + 886033068781154065770*a(n-26) - 4721965559520496562627*a(n-27) - 5509376950388514662900*a(n-28) + 10653562229545596548524*a(n-29) + 18834020448070528000425*a(n-30) - 13713203739273583437918*a(n-31) - 39901048103937802002297*a(n-32) + 5990235505236935045275*a(n-33) + 55578455940464802386193*a(n-34) + 9833720985091163714610*a(n-35) - 53199364735680239822375*a(n-36) - 20471708203330957085189*a(n-37) + 36360819633145359083726*a(n-38) + 19039459457994629613832*a(n-39) - 18441722899588867748260*a(n-40) - 11149442822036811116540*a(n-41) + 7236658075872597214567*a(n-42) + 4449967250415200774117*a(n-43) - 2282182298554463978746*a(n-44) - 1228251073290200895900*a(n-45) + 585623283999198391080*a(n-46) + 226322307178914076710*a(n-47) - 118031343431948320300*a(n-48) - 24333902651539947292*a(n-49) + 17131550746937484216*a(n-50) + 699869866498272640*a(n-51) - 1545249253168883776*a(n-52) + 148360308338197984*a(n-53) + 62588687728118848*a(n-54) - 14435746763377920*a(n-55) + 239262718302208*a(n-56) + 190949696145408*a(n-57) - 13490873376768*a(n-58) - 661368913920*a(n-59) + 68133519360*a(n-60)
\end{quote}
As of the ``Last modified'' line on the live entry (Jun 02 2025 09:14:48, revision 7) the statement is
still recorded as empirical, and nothing on the entry records it as settled either way.

Written out, the claim is
\begin{equation}\label{eq:conj}
a(n)\;=\;149\,a(n-1) - 8814\,a(n-2) + 256435\,a(n-3) - 3393920\,a(n-4) + 449245\,a(n-5) + 514842621\,a(n-6) - 4157999842\,a(n-7) - 23261708552\,a(n-8) + 407861534972\,a(n-9) + 96979298270\,a(n-10) - 20363478746809\,a(n-11) + 30533472981167\,a(n-12) + 660568556988005\,a(n-13) - 1408051468806556\,a(n-14) - 15383126507976075\,a(n-15) + 32902261921526205\,a(n-16) + 269077534259492089\,a(n-17) - 467962804548993846\,a(n-18) - 3572775805160926106\,a(n-19) + 4038250887549205858\,a(n-20) + 35582810698357101463\,a(n-21) - 17005381783517699700\,a(n-22) - 259443464221965708471\,a(n-23) - 32542572652022730119\,a(n-24) + 1341569768136539281807\,a(n-25) + 886033068781154065770\,a(n-26) - 4721965559520496562627\,a(n-27) - 5509376950388514662900\,a(n-28) + 10653562229545596548524\,a(n-29) + 18834020448070528000425\,a(n-30) - 13713203739273583437918\,a(n-31) - 39901048103937802002297\,a(n-32) + 5990235505236935045275\,a(n-33) + 55578455940464802386193\,a(n-34) + 9833720985091163714610\,a(n-35) - 53199364735680239822375\,a(n-36) - 20471708203330957085189\,a(n-37) + 36360819633145359083726\,a(n-38) + 19039459457994629613832\,a(n-39) - 18441722899588867748260\,a(n-40) - 11149442822036811116540\,a(n-41) + 7236658075872597214567\,a(n-42) + 4449967250415200774117\,a(n-43) - 2282182298554463978746\,a(n-44) - 1228251073290200895900\,a(n-45) + 585623283999198391080\,a(n-46) + 226322307178914076710\,a(n-47) - 118031343431948320300\,a(n-48) - 24333902651539947292\,a(n-49) + 17131550746937484216\,a(n-50) + 699869866498272640\,a(n-51) - 1545249253168883776\,a(n-52) + 148360308338197984\,a(n-53) + 62588687728118848\,a(n-54) - 14435746763377920\,a(n-55) + 239262718302208\,a(n-56) + 190949696145408\,a(n-57) - 13490873376768\,a(n-58) - 661368913920\,a(n-59) + 68133519360\,a(n-60)
\end{equation}
for all $n>60$.

\section{The condition is local, so the rows form a digraph}

Write an array of the shape in question as its list of rows
$r^{(1)},r^{(2)},\dots$, each an element of
$V=\{0,1,\dots,5\}^{5}$, so $|V|=S=7776$. A $2\times2$ subblock of the array
occupies two consecutive rows and two consecutive columns; if the two rows are $r$
and $s$ (in that order) and the two columns are numbered $j,j+1$,
the subblock is
\[
\begin{pmatrix}\alpha&\beta\\\gamma&\delta\end{pmatrix}=\begin{pmatrix}r_j&r_{j+1}\\s_j&s_{j+1}\end{pmatrix}.
\]
The entry's requirement on that subblock is
\begin{equation}\label{eq:loc}
\max\bigl(|\beta-\alpha|,|\delta-\beta|,|\gamma-\delta|,|\alpha-\gamma|,|\delta-\alpha|,|\gamma-\beta|\bigr)\;\le\;1,
\end{equation}
a condition on the four entries alone.

For $r,s\in V$ declare $r\to s$ an edge of a digraph $G$ when, for every
$1\le j\le 5-1$,
\begin{itemize}
\item the four entries above satisfy \eqref{eq:loc};

\end{itemize}
Let $M\in\{0,1\}^{S\times S}$ be the adjacency matrix of $G$ and $\mathbf{1}$ the
all-ones vector.

\begin{lemma}\label{lem:walk}
The arrays counted by the entry with $L$ rows are exactly the walks
$r^{(1)}\to r^{(2)}\to\cdots\to r^{(L)}$ of length $L-1$ in $G$. Consequently
\[
a(n)\;=\;\mathbf{1}^{\!\top}M^{n}\mathbf{1} .
\]
\end{lemma}

\begin{proof}
An array with $L$ rows and $5$ columns is precisely a sequence
$r^{(1)},\dots,r^{(L)}$ of elements of $V$; conversely any such sequence is an array.
Every $2\times2$ subblock lies in two consecutive rows $r^{(t)},r^{(t+1)}$ and two
consecutive columns, so the array satisfies the entry's condition on all of its subblocks
if and only if each consecutive pair $\bigl(r^{(t)},r^{(t+1)}\bigr)$ satisfies it for
every column pair --- which is exactly the edge relation of $G$. The number of walks of
length $L-1$, summed over all starting and ending vertices, is
$\mathbf{1}^{\!\top}M^{L-1}\mathbf{1}$, since $(M^{k})_{r,s}$ counts the walks of
length $k$ from $r$ to $s$. The entry's index $n$ corresponds to $L=n+1$ rows.
\end{proof}

\section{The criterion}

Let $q(t)=t^{60}-\sum_i c_i\,t^{60-i}$ be the characteristic polynomial of
\eqref{eq:conj}, with the $c_i$ read off that equation.

\begin{lemma}\label{lem:crit}
For every $n$ with $n+1-1\ge60$,
\[
a(n)-\sum_i c_i\,a(n-i)\;=\;\mathbf{1}^{\!\top}M^{\,n+1-1-60}\,q(M)\,\mathbf{1} .
\]
Hence \eqref{eq:conj} holds from that point on if and only if
$\mathbf{1}^{\!\top}M^{j}q(M)\mathbf{1}=0$ for every $j\ge0$; and it is enough to check
$j=0,1,\dots,S-1$.
\end{lemma}

\begin{proof}
By Lemma~\ref{lem:walk}, $a(n-i)=\mathbf{1}^{\!\top}M^{\,n+1-1-i}\mathbf{1}$, so
\[
a(n)-\sum_i c_i a(n-i)
=\mathbf{1}^{\!\top}\Bigl(M^{m}-\sum_i c_i M^{m-i}\Bigr)\mathbf{1}
=\mathbf{1}^{\!\top}M^{\,m-60}\,q(M)\,\mathbf{1}
\]
with $m=n+1-1$, which is the stated identity; the equivalence follows by letting
$m-60=j$ run over $\ge0$. For the last clause put $w=q(M)\mathbf{1}$. By the
Cayley--Hamilton theorem $M^{S}$ is an integer combination of $I,M,\dots,M^{S-1}$, so
$\mathbf{1}^{\!\top}M^{j}w$ for $j\ge S$ is the corresponding combination of the earlier
values; if those all vanish, so do all the rest.
\end{proof}

\section{The computation}

\begin{theorem}
The recurrence \eqref{eq:conj} holds for every $n>60$, which is the statement of
Section~1.
\end{theorem}

\begin{proof}
The digraph was constructed from \eqref{eq:loc} by a depth-first scan across the $5$
columns: the edge condition is a conjunction of one constraint per consecutive column
pair, so the successors of a row $r$ are enumerated column by column without testing all
$S^2$ pairs. The vector $w=q(M)\mathbf{1}\in\mathbb{Z}^{S}$ was then formed by $60$
matrix--vector products in exact integer arithmetic, and $\mathbf{1}^{\!\top}M^{j}w$ was
evaluated for $j=0,1,\dots$, again exactly. Every one of those integers vanishes from the
index corresponding to $n=60+1$ onwards, and the run continued until $w$ itself was the
zero vector, which settles every larger $j$ at once. By Lemma~\ref{lem:crit} the recurrence
holds for every $n>60$.
\end{proof}

\begin{remark}
The argument gives more than the recurrence: it shows the sequence is $C$-finite with
characteristic polynomial dividing that of $M$, so it has a rational generating function
whose denominator divides $\det(I-xM)$. The conjectured recurrence is one consequence; the
minimal one is a divisor of $q$.
\end{remark}

\section{Verification}

Three checks, each independent of the algebra above.

First, the digraph was built directly from the entry's stated condition and the walk counts
$\mathbf{1}^{\!\top}M^{L-1}\mathbf{1}$ compared against the entry's DATA at
every one of the $13$ published indices; the values agree exactly. That is what ties
the digraph to the sequence: a model that reproduced only some of the published terms would
be the wrong model, and would have been rejected. In particular it is what rules out a
misreading of the entry's wording, since a different reading of the condition gives a
different digraph and different counts.

Second, the recurrence \eqref{eq:conj} was evaluated on those published terms in exact
integer arithmetic, with no matrices involved, and vanishes wherever it applies.

Third, the annihilation test of Lemma~\ref{lem:crit} was run in exact integer arithmetic
until the working vector was identically zero, not to a sampled prefix, and returned zero
throughout.

\begin{thebibliography}{9}
\bibitem{oeis} The OEIS Foundation, \emph{The On-Line Encyclopedia of Integer
Sequences}, \url{https://oeis.org}, 2026. Sequence A234180.
\bibitem{stanley} R.~P.~Stanley, \emph{Enumerative Combinatorics}, Volume 1, 2nd ed.,
Cambridge University Press, 2011. (Transfer-matrix method, Section 4.7.)
\bibitem{fs} P.~Flajolet and R.~Sedgewick, \emph{Analytic Combinatorics}, Cambridge
University Press, 2009.
\bibitem{horn} R.~A.~Horn and C.~R.~Johnson, \emph{Matrix Analysis}, 2nd ed., Cambridge
University Press, 2013.
\end{thebibliography}

\end{document}
